\begin{enumerate}
		\item \textbf{Affirmation 1} :

		$f$ est dérivable  et, pour tout $x$ réel : \quad $f(x) = (x - 2)\e^x$

		donc, pour tout $x$ réel :\quad $f'(x)=1\times \e^x + (x-2)\times \e^x = ( x-1)\e^x$.

		Ainsi, pour tout réel $x$,\quad $\begin{aligned}[t]
		f'(x)-2f(x) & =(x-1)\e^x-2(x-2)\e^x \\
		& =(-x+3)\e^x \\
		& = -(x-3)\e^x
		\end{aligned}$

		L'affirmation est \textbf{fausse}.

		\item \textbf{Affirmation 2} :

		On pose, pour tout $x$ réel, \quad $v(x)=2x+1$ \quad et \quad $u'(x)=\cos x$.

		On a alors, pour tout $x$ réel :\quad $v'(x)=2$ \quad  et on peut choisir \quad $u(x)=\sin x$.

		Par intégration par parties, on a :\quad
		$\begin{aligned}[t]
		\displaystyle\int_0^{\pi} f(x)\:\text{d}x &= \bigg[ (2x+1) \sin x \bigg]_0^{\pi} - \displaystyle \int_0^{\pi} 2 \sin x \d x\\
		&= \bigg( (2 \pi +1) \sin \pi - \sin 0 \bigg) - \bigg[-2 \cos x \bigg]_0^{\pi} \\
		&= 0 + 2 \cos \pi -2 \cos 0 \\
		&= -4
		\end{aligned}$

		L'affirmation est \textbf{fausse}.

		\item \textbf{Affirmation 3} :

		\begin{itemize}
			\item La fonction du second degré \quad $x \longmapsto 2x^2-6x-56$\quad admet deux racines :\quad $-4$ et 7.

			Le coefficient du terme de degré 2 est \quad $a=2$\quad  donc cette fonction est strictement positive sur $]7 ; +\infty[$.

			\item La fonction affine \quad $x \longmapsto 4x-8$\quad  est strictement positive sur $]2 ;+\infty[$.
		\end{itemize}

		On peut résoudre l'inéquation sur l'intervalle $]7 ; +\infty[$, et on a :

		$\begin{aligned}[t] \forall x \in ]7 ;+\infty[, \quad
		\ln \left(2x^2 - 6x - 56\right) \geqslant \ln(4x - 8)
		& \iff 2x^2 - 6x - 56 \geqslant 4x - 8 \quad &\text{car la fonction ln est}\\
		&&\text{croissante sur }\R^{*+}\\
		& \iff 2x^2 - 10x - 48 \geqslant 0 \\
		& \iff x^2-5x-24 \geqslant 0&\text{car }2 >0
		\end{aligned}$

		Le discriminant est :\quad $\Delta = b^2-4ac = 121$,\quad donc la fonction du second degré \quad $x \longmapsto x^2-5x-24$ \quad admet deux racines :\quad $-3$ et 8.

		Le coefficient du terme de degré 2 est $a=1$, donc cette fonction est positive sur $]-\infty  ; -3] \cup [8 ;+\infty[$.

		Sur $]7  ; + \infty[$, l'inéquation a pour ensemble de solutions l'intervalle $[8 ; +\infty[$.

		L'affirmation est \textbf{vraie}.

		\item \textbf{Affirmation 4} :

		Soit $n$ un entier naturel :

		$v_{n+1} = \dfrac{2}{u_{n+1}} = \dfrac{2}{~\dfrac{3u_n}{5u_n+3}~} = 2\times \dfrac{5u_n + 3}{3u_n} = 2\times \left(\dfrac{5u_n}{3u_n} + \dfrac{3}{3u_n}\right) = 2\times \left(\dfrac{5}{3} + \dfrac{1}{u_n}\right) = \dfrac{10}{3} + \dfrac{2}{u_n} = \dfrac{10}{3} + v_n$.

		On reconnaît la relation de récurrence d'une suite arithmétique, de raison $r = \dfrac{10}{3}$.

		L'affirmation est \textbf{vraie}.

		\item \textbf{Affirmation 5} :

		On considère l'expérience aléatoire "obtenir un 6 en lançant un dé équilibré à 6 face numéroté de 1 à 6" a deux issues : le succès de probabilité $\dfrac16$ et l'échec de probabilité $\dfrac56$. Il s'agit donc d'une épreuve de Bernoulli.

		La variable aléatoire qui compte le nombre de succès (les apparitions du 6) quand on répète \np{3600} fois cette épreuve suit la loi binomiale de paramètre $p=\dfrac16$ et $n=\np{3600}$.

		On a alors \quad $E(X)=\dfrac16 \times 3600 = 600$, \quad et \quad $V(X)=\dfrac16 \times \dfrac56 \times 3600=500$.

		D'après l'inégalité de Bienaymé-Tchebytchev, pour tout $t >0$, on a :\quad $P \bigg( \big\vert X-600 \big\vert \geqslant t \bigg) \leqslant \dfrac{500}{t^2}$.

		Or \quad $P \bigg( 450 < X < 750 \bigg) = P \bigg( \big\vert X-600 \big\vert < 150 \bigg) =1- P \bigg( \big\vert X-600 \big\vert \geqslant 150 \bigg)$.

		En appliquant l'inégalité de Bienaymé-Tchebytchev pour $t=150$, on obtient donc :

		$P \bigg( \big\vert X-600 \big\vert \geqslant 150 \bigg) \leqslant \dfrac{500}{150^2}$

		On en déduit que \quad $P \bigg( 450 < X < 750 \bigg)  \geqslant 1-\dfrac{500}{150^2}$.

		Or, \quad $1-\dfrac{500}{150^2} \approx 0,978$.

		On peut donc conclure : \quad $P \bigg( 450 < X < 750 \bigg) \geqslant \np{0,97}$.

		L'affirmation est \textbf{vraie}.
	\end{enumerate}
